\section{实验总结与展望}\label{sec:conclusion}

本报告基于 162 张带像素级标签的人脸图像，完整呈现了从实验三（YCrCb 颜色建模与基础连通域后处理）到实验四（双眼定位引导的双椭圆空间约束与低支持度放宽阈值重建）的递进研究过程。全文的核心工作与结论可总结为以下四个方面：

第一，\textbf{YCrCb 色度建模与矩形阈值的覆盖特性：}训练集 638.8 万个皮肤与鼻子像素在 Cr/Cb 平面上呈明显的倾斜椭圆聚集（两通道标准差均约为 9.4，相关系数约为 $-0.7644$）。实验三选定的 $k=1.25$ 矩形阈值侧重于抑制非肤色背景，实质上截取了高密度肤色核心区；这一特性保证了常规光照下的高精确率，但在暗光或冷偏色图像上容易导致面部候选碎裂。

第二，\textbf{基础连通域后处理的阶段增益：}在实验三中，$k=1.25$ 阈值配合最大 8 连通域筛选与内部空洞填充，在 32 张测试图上取得全局 F1 70.69\%、IoU 54.67\%，比纯阈值基线分别提升 8.03 和 9.04 个百分点。分阶段像素分解表明，阈值收紧、最大连通域筛选与内部填孔的正负像素交换率分别为 $1:1.6$、$1:8.3$ 和 $3.3:1$，三者协同完成了背景去噪与封闭五官恢复。

第三，\textbf{纯颜色与无边界形态学操作的等值线瓶颈：}误差语义分解显示，头发（60.58\%）与脖颈（23.40\%）合计占实验三最终误检的 83.98\%。在缺乏人脸几何边界的情况下，$41\times41$ 椭圆核“膨胀—填孔—腐蚀”补全虽将测试集召回率提高 4.00 个百分点（达 79.55\%），但新增前景像素中 64.55\% 属于外围头发与背景误检，精确率下降 2.80 个百分点，全局 F1（70.70\%）保持不变。这证明无空间锚点的形态学操作只能在 F1 等值线上调节精确率与召回率的分配。

第四，\textbf{双眼引导双椭圆约束与低支持重建的协同突破：}实验四通过分阶段连通域诊断与轮廓直接拟合实验，证实了待修整肤色轮廓本身易受粘连头发或局部碎片带偏；进而利用灰度黑帽变换结合训练集左右眼独立 $3\sigma$ 窗口先验定位双眼（验证集 $e\le 0.25$ 准确率达 28/32，配合严格/放宽两级面积比回退覆盖 31/32 张图，无配对样本自动回退 Lab3），在双眼局部坐标系下生成正中及左右侧向候选（$\pm 0.16\,d_{\text{eye}}$，偏移门槛 $\Delta S_{\text{lat}}=0.06$）与内外嵌套双椭圆（$\lambda_{\text{in}}=0.82,\;\lambda_{\text{out}}=1.20$）。外椭圆有效截断了大面积粘连的头发与脖颈误检，内椭圆针对支持率 $\rho_{\text{in}}<0.18$ 的碎裂样本触发 $k_{\text{rec}}=2.00$ 放宽肤色重选与 $27\times27$ 闭运算受限重建，最后经内椭圆 $11\times11$ 强闭运算与过渡带 $5\times5$ 弱闭运算完成核心种子锚定连接。在 32 张验证图上，实验四最终方案将全局 Precision 从 64.86\% 提升至 79.88\%（+15.02 个百分点），Recall 从 82.70\% 提升至 86.15\%（+3.45 个百分点），全局 F1 从 72.70\% 跃升至 82.89\%（+10.19 个百分点），IoU 从 57.11\% 提升至 70.79\%（+13.68 个百分点），实现了精确率与召回率的同步跃升。

展望后续工作，可在现有颜色—几何混合框架下沿三个方向进一步拓展：一是将 Cr/Cb 轴对齐矩形阈值升级为考虑协方差结构的马氏距离椭圆判别；二是结合鼻部与嘴部几何线索优化侧向椭圆姿态选择准则，充分释放侧脸修正潜力（理想姿态上界 F1 达 84.44\%）；三是在内椭圆上缘引入局部边缘梯度与纹理约束，进一步抑制暖调刘海在放宽阈值重建时的误检混入。
